\documentclass[12pt,a4paper]{article}
\usepackage[utf8]{vietnam}
\usepackage{amsmath,amssymb}
\usepackage{tabularx}
\usepackage{tikz}
\usepackage{xcolor}
% Giảm lề tối đa tận dụng diện tích giấy (chuẩn máy in laser/photo)
\usepackage[top=0.6cm, bottom=0.6cm, left=0.6cm, right=0.6cm]{geometry}
\usepackage[dethi]{ex_test}

\setlength{\parskip}{0pt}

% Dòng kẻ chấm mờ đều đặn chuẩn viết tay (khoảng cách 7.2mm rộng rãi)
\newcommand{\dongke}[1]{%
  \begingroup%
  \setlength{\parskip}{0pt}%
  \baselineskip=20.5pt%
  \vspace{1.5mm}%
  \foreach \i in {1,...,#1}{%
    \noindent{\color{gray!60}\dotfill}\par%
  }%
  \vspace{1mm}%
  \endgroup%
}

% Định nghĩa ô bài tập tự luận trong bảng
\newcommand{\oBaiTap}[3][7]{%
  \vspace{1.5mm}%
  \noindent\textbf{#2} #3\par%
  \dongke{#1}%
}

% Câu trắc nghiệm 1 cột: 4 đáp án dàn hàng ngang trên 1 dòng
\newcommand{\cauTN}[6]{% \cauTN{số câu}{đề}{A}{B}{C}{D}
  \noindent\textbf{Câu #1.} #2\par\vspace{0.4mm}%
  \noindent%
  \makebox[0.25\linewidth][l]{\textbf{A.} #3}%
  \makebox[0.25\linewidth][l]{\textbf{B.} #4}%
  \makebox[0.25\linewidth][l]{\textbf{C.} #5}%
  \makebox[0.25\linewidth][l]{\textbf{D.} #6}\par\vspace{0.8mm}%
}

\begin{document}
\pagestyle{empty}

% #############################################################################
% ========================= TOÁN 6 - TUẦN 2 - BUỔI 1: MÃ ĐỀ 101 =============
% #############################################################################

% -------------------- TRANG 1 (MÃ ĐỀ 101): TRẮC NGHIỆM --------------------
\noindent
\begin{tabularx}{\linewidth}{|X|c|}
\hline
\textbf{BÀI 3: THỨ TỰ SỐ TỰ NHIÊN (TUẦN 2 - BUỔI 1)} & \textbf{Mã đề: 101} \\
\hline
\end{tabularx}

\vspace{1.2mm}

\noindent
\begin{tabularx}{\linewidth}{|X|}
\hline
\rule[-3mm]{0pt}{9.5mm}\textbf{Họ và tên:} \dotfill \quad \textbf{Lớp:} \makebox[1.8cm]{\dotfill} \\
\hline
\end{tabularx}

\vspace{1.2mm}

\noindent
\begin{tabularx}{\linewidth}{|p{3cm}|X|}
\hline
\centering\textbf{ĐIỂM SỐ} & \centering\textbf{LỜI NHẬN XÉT} \tabularnewline
\hline
\rule{0pt}{2.2cm} & \tabularnewline
\hline
\end{tabularx}

\vspace{1.5mm}
\noindent\textbf{BẢNG ĐIỀN ĐÁP ÁN TRẮC NGHIỆM (Mỗi câu đúng 0,25 điểm):}
\vspace{1.2mm}

\noindent
\renewcommand{\arraystretch}{1.5}
\begin{tabularx}{\linewidth}{|c|*{12}{>{\centering\arraybackslash}X|}}
\hline
\textbf{Câu} & \textbf{1} & \textbf{2} & \textbf{3} & \textbf{4} & \textbf{5} & \textbf{6} & \textbf{7} & \textbf{8} & \textbf{9} & \textbf{10} & \textbf{11} & \textbf{12} \\
\hline
\textbf{Đ/A} & & & & & & & & & & & & \\
\hline
\end{tabularx}
\renewcommand{\arraystretch}{1.0}

\vspace{3mm}

\noindent\textbf{\large PHẦN I. TRẮC NGHIỆM KHÁCH QUAN (3,0 điểm)}
\vspace{1.5mm}

\cauTN{1}{Trên tia số nằm ngang, nếu $a < b$ thì điểm $a$ nằm ở:}{Bên phải điểm $b$}{Bên trái điểm $b$}{Trùng điểm $b$}{Phía trên điểm $b$}

\cauTN{2}{Khẳng định nào sau đây là \textbf{sai}?}{$2\,024 > 2\,023$}{$10\,000 > 9\,999$}{$54\,321 < 54\,312$}{$102 < 120$}

\cauTN{3}{Trong các số $4\,567; 4\,576; 4\,765; 4\,675$, số lớn nhất là:}{$4\,567$}{$4\,576$}{$4\,765$}{$4\,675$}

\cauTN{4}{Kí hiệu $x \le 6$ có nghĩa là:}{$x < 6$}{$x = 6$}{$x < 6$ hoặc $x = 6$}{$x > 6$ hoặc $x = 6$}

\cauTN{5}{Tập hợp các số tự nhiên $x$ thoả mãn $3 < x \le 7$ là:}{$\{3; 4; 5; 6\}$}{$\{4; 5; 6; 7\}$}{$\{3; 4; 5; 6; 7\}$}{$\{4; 5; 6\}$}

\cauTN{6}{Chữ số thích hợp điền vào dấu $*$ để $3*5 < 315$ là:}{$0$}{$1$}{$2$}{$3$}

\cauTN{7}{Cho biết $m > n$ và $n > p$. Khẳng định nào chắc chắn đúng?}{$m < p$}{$m > p$}{$m = p$}{$n < p$}

\cauTN{8}{Số tự nhiên liền sau số lớn nhất có ba chữ số là:}{$998$}{$1\,000$}{$999$}{$1\,001$}

\cauTN{9}{Số liền trước của số tự nhiên $a + 2$ (với $a \in \mathbb{N}$) là:}{$a + 3$}{$a + 1$}{$a$}{$a - 1$}

\cauTN{10}{Tập hợp $A = \{x \in \mathbb{N} \mid 15 \le x \le 25\}$ có số phần tử là:}{$10$}{$11$}{$9$}{$12$}

\cauTN{11}{Sắp xếp các số $12; 0; 8; 25; 19$ theo thứ tự giảm dần là:}{$0; 8; 12; 19; 25$}{$25; 19; 12; 8; 0$}{$25; 19; 8; 12; 0$}{$0; 25; 19; 12; 8$}

\cauTN{12}{Có bao nhiêu số tự nhiên $x$ thỏa mãn $198 \le x < 203$?}{$4$}{$5$}{$6$}{$7$}

\newpage
% -------------------- TRANG 2 (MÃ ĐỀ 101): TỰ LUẬN --------------------
\noindent
\begin{tabularx}{\linewidth}{|X|c|}
\hline
\textbf{PHẦN II. TỰ LUẬN (7,0 ĐIỂM)} & \textbf{Mã đề: 101} \\
\hline
\end{tabularx}

\vspace{1.2mm}

\noindent
\begin{tabularx}{\linewidth}{|X|}
\hline
\rule[-3mm]{0pt}{9.5mm}\textbf{Họ và tên:} \dotfill \quad \textbf{Lớp:} \makebox[1.8cm]{\dotfill} \\
\hline
\end{tabularx}

\vspace{2.5mm}

\noindent\textbf{\large Câu 1 (2,5 điểm). So sánh và Viết tập hợp:}
\vspace{1.2mm}

\noindent
\begin{tabularx}{\linewidth}{|X|X|}
\hline
\oBaiTap[7]{a)}{Cho các số sau: $3\,205; 3\,052; 3\,520; 3\,025; 3\,250$. Hãy sắp xếp các số đã cho theo thứ tự từ bé đến lớn.} & 
\oBaiTap[7]{b)}{Viết tập hợp $M = \{x \in \mathbb{N} \mid 14 \le x < 20\}$ bằng cách liệt kê các phần tử. Tập hợp $M$ có bao nhiêu phần tử?} \\
\hline
\end{tabularx}

\vspace{3.5mm}

\noindent\textbf{\large Câu 2 (2,5 điểm). Điền chữ số và Tìm số tự nhiên:}
\vspace{1.2mm}

\noindent
\begin{tabularx}{\linewidth}{|X|X|}
\hline
\oBaiTap[7]{a)}{Tìm chữ số thích hợp thay cho dấu $*$ để được khẳng định đúng: $5*3 > 583$ và $43* < 432$.} & 
\oBaiTap[7]{b)}{Tìm số tự nhiên có hai chữ số $\overline{ab}$, biết rằng $20 < \overline{ab} < 35$ và chữ số hàng đơn vị lớn hơn chữ số hàng chục $3$ đơn vị.} \\
\hline
\end{tabularx}

\vspace{3.5mm}

\noindent\textbf{\large Câu 3 (2,0 điểm). Bài toán thực tế:}
\vspace{1.2mm}

\noindent
\begin{tabularx}{\linewidth}{|X|}
\hline
\vspace{0.8mm}
Bảng sau cho biết độ cao của các đỉnh núi: Fansipan ($3\,143\text{ m}$), Pu Si Lung ($3\,083\text{ m}$), Pu Ta Leng ($3\,049\text{ m}$), Bạch Mộc Lương Tử ($3\,046\text{ m}$).\\
a) Sắp xếp tên các đỉnh núi theo thứ tự độ cao giảm dần.\\
b) Tính độ chênh lệch chiều cao giữa đỉnh núi cao nhất và thấp nhất trong các đỉnh núi trên.
\par\vspace{1.5mm}
\noindent\textit{Bài giải:}
\dongke{7}
\\ \hline
\end{tabularx}

\newpage
% #############################################################################
% ========================= TOÁN 6 - TUẦN 2 - BUỔI 1: MÃ ĐỀ 102 =============
% #############################################################################

% -------------------- TRANG 3 (MÃ ĐỀ 102): TRẮC NGHIỆM --------------------
\noindent
\begin{tabularx}{\linewidth}{|X|c|}
\hline
\textbf{BÀI 3: THỨ TỰ SỐ TỰ NHIÊN (TUẦN 2 - BUỔI 1)} & \textbf{Mã đề: 102} \\
\hline
\end{tabularx}

\vspace{1.2mm}

\noindent
\begin{tabularx}{\linewidth}{|X|}
\hline
\rule[-3mm]{0pt}{9.5mm}\textbf{Họ và tên:} \dotfill \quad \textbf{Lớp:} \makebox[1.8cm]{\dotfill} \\
\hline
\end{tabularx}

\vspace{1.2mm}

\noindent
\begin{tabularx}{\linewidth}{|p{3cm}|X|}
\hline
\centering\textbf{ĐIỂM SỐ} & \centering\textbf{LỜI NHẬN XÉT} \tabularnewline
\hline
\rule{0pt}{2.2cm} & \tabularnewline
\hline
\end{tabularx}

\vspace{1.5mm}
\noindent\textbf{BẢNG ĐIỀN ĐÁP ÁN TRẮC NGHIỆM (Mỗi câu đúng 0,25 điểm):}
\vspace{1.2mm}

\noindent
\renewcommand{\arraystretch}{1.5}
\begin{tabularx}{\linewidth}{|c|*{12}{>{\centering\arraybackslash}X|}}
\hline
\textbf{Câu} & \textbf{1} & \textbf{2} & \textbf{3} & \textbf{4} & \textbf{5} & \textbf{6} & \textbf{7} & \textbf{8} & \textbf{9} & \textbf{10} & \textbf{11} & \textbf{12} \\
\hline
\textbf{Đ/A} & & & & & & & & & & & & \\
\hline
\end{tabularx}
\renewcommand{\arraystretch}{1.0}

\vspace{3mm}

\noindent\textbf{\large PHẦN I. TRẮC NGHIỆM KHÁCH QUAN (3,0 điểm)}
\vspace{1.5mm}

\cauTN{1}{Trên tia số nằm ngang, nếu $x > y$ thì điểm $x$ nằm ở:}{Bên trái điểm $y$}{Bên phải điểm $y$}{Trùng điểm $y$}{Phía dưới điểm $y$}

\cauTN{2}{Khẳng định nào sau đây là \textbf{sai}?}{$3\,025 > 3\,024$}{$9\,998 < 10\,000$}{$76\,543 < 76\,534$}{$230 > 203$}

\cauTN{3}{Trong các số $6\,789; 6\,879; 6\,798; 6\,897$, số nhỏ nhất là:}{$6\,789$}{$6\,879$}{$6\,798$}{$6\,897$}

\cauTN{4}{Kí hiệu $a \ge 8$ có nghĩa là:}{$a > 8$}{$a = 8$}{$a > 8$ hoặc $a = 8$}{$a < 8$ hoặc $a = 8$}

\cauTN{5}{Tập hợp các số tự nhiên $y$ thoả mãn $5 \le y < 9$ là:}{$\{5; 6; 7; 8\}$}{$\{6; 7; 8; 9\}$}{$\{5; 6; 7; 8; 9\}$}{$\{6; 7; 8\}$}

\cauTN{6}{Chữ số thích hợp điền vào dấu $*$ để $7*2 > 782$ là:}{$7$}{$8$}{$9$}{$6$}

\cauTN{7}{Cho biết $a < b$ và $b < c$. Khẳng định nào chắc chắn đúng?}{$a > c$}{$a < c$}{$a = c$}{$b > c$}

\cauTN{8}{Số tự nhiên liền trước số nhỏ nhất có bốn chữ số là:}{$999$}{$1\,001$}{$990$}{$1\,000$}

\cauTN{9}{Số liền sau của số tự nhiên $x - 1$ (với $x \in \mathbb{N}^*$) là:}{$x - 2$}{$x$}{$x + 1$}{$x + 2$}

\cauTN{10}{Tập hợp $B = \{x \in \mathbb{N} \mid 20 \le x \le 32\}$ có số phần tử là:}{$12$}{$13$}{$11$}{$14$}

\cauTN{11}{Sắp xếp các số $15; 3; 0; 28; 21$ theo thứ tự tăng dần là:}{$28; 21; 15; 3; 0$}{$0; 3; 15; 21; 28$}{$0; 3; 15; 28; 21$}{$3; 0; 15; 21; 28$}

\cauTN{12}{Có bao nhiêu số tự nhiên $x$ thỏa mãn $299 < x \le 305$?}{$5$}{$6$}{$7$}{$4$}

\newpage
% -------------------- TRANG 4 (MÃ ĐỀ 102): TỰ LUẬN --------------------
\noindent
\begin{tabularx}{\linewidth}{|X|c|}
\hline
\textbf{PHẦN II. TỰ LUẬN (7,0 ĐIỂM)} & \textbf{Mã đề: 102} \\
\hline
\end{tabularx}

\vspace{1.2mm}

\noindent
\begin{tabularx}{\linewidth}{|X|}
\hline
\rule[-3mm]{0pt}{9.5mm}\textbf{Họ và tên:} \dotfill \quad \textbf{Lớp:} \makebox[1.8cm]{\dotfill} \\
\hline
\end{tabularx}

\vspace{2.5mm}

\noindent\textbf{\large Câu 1 (2,5 điểm). So sánh và Viết tập hợp:}
\vspace{1.2mm}

\noindent
\begin{tabularx}{\linewidth}{|X|X|}
\hline
\oBaiTap[7]{a)}{Cho các số sau: $4\,306; 4\,063; 4\,630; 4\,036; 4\,360$. Hãy sắp xếp các số đã cho theo thứ tự từ bé đến lớn.} & 
\oBaiTap[7]{b)}{Viết tập hợp $N = \{x \in \mathbb{N} \mid 23 < x \le 29\}$ bằng cách liệt kê các phần tử. Tập hợp $N$ có bao nhiêu phần tử?} \\
\hline
\end{tabularx}

\vspace{3.5mm}

\noindent\textbf{\large Câu 2 (2,5 điểm). Điền chữ số và Tìm số tự nhiên:}
\vspace{1.2mm}

\noindent
\begin{tabularx}{\linewidth}{|X|X|}
\hline
\oBaiTap[7]{a)}{Tìm chữ số thích hợp thay cho dấu $*$ để được khẳng định đúng: $6*7 < 617$ và $85* > 857$.} & 
\oBaiTap[7]{b)}{Tìm số tự nhiên có hai chữ số $\overline{xy}$, biết rằng $30 < \overline{xy} < 45$ và chữ số hàng đơn vị kém chữ số hàng chục $2$ đơn vị.} \\
\hline
\end{tabularx}

\vspace{3.5mm}

\noindent\textbf{\large Câu 3 (2,0 điểm). Bài toán thực tế:}
\vspace{1.2mm}

\noindent
\begin{tabularx}{\linewidth}{|X|}
\hline
\vspace{0.8mm}
Bảng sau cho biết diện tích các địa phương: Hà Nội ($3\,359\text{ km}^2$), TP. Hồ Chí Minh ($2\,095\text{ km}^2$), Đà Nẵng ($1\,285\text{ km}^2$), Hải Phòng ($1\,562\text{ km}^2$).\\
a) Sắp xếp tên các tỉnh/thành phố theo thứ tự diện tích giảm dần.\\
b) Tính độ chênh lệch diện tích giữa địa phương lớn nhất và nhỏ nhất trong các địa phương trên.
\par\vspace{1.5mm}
\noindent\textit{Bài giải:}
\dongke{7}
\\ \hline
\end{tabularx}

\end{document}
